\section{Introduction}
\label{sec:intro}

Hardware specialization enables gains in performance and energy efficiency through architectures such as GPUs and tensor accelerators~\cite{autocomp}. Exploiting these capabilities has traditionally required substantial human effort: experts must organize data movement, select instructions, and tune kernels for the target hardware. Each new architecture can require developers to revisit these choices and adapt existing implementations to a different programming interface.

CPU SIMD and vector units are a widely used form of this specialization. The RISC-V Vector extension (RVV) supports vector processing across implementations with different hardware vector lengths~\cite{rvv}. Building library support for RVV makes these capabilities available to applications and has motivated work on translating existing vector code~\cite{intrintrans,saltyrn}. Libraries with NEON or AVX implementations already provide optimized kernels that can serve as the source for these translations. The execution models differ, however: NEON processes vectors of a fixed width, whereas RVV selects the number of active elements at runtime~\cite{acle,rvv}. Both expose vector operations through C intrinsic interfaces, allowing programmers to combine ordinary C control flow with operations that the compiler lowers to hardware instructions.

Recent work has shown that agents can generate and optimize kernels using hardware knowledge and execution feedback~\cite{kernelbench,autocomp,saltyrn}. For SIMD code, the C intrinsic interface lets an agent apply this knowledge directly to instruction selection and loop organization. However, agents may exploit gaps in test coverage through reward hacking, producing kernels that pass the available tests without preserving the intended behavior~\cite{kernelbench-eval}. Existing work addresses correctness through testing and formal validation, including bounded equivalence checking of translated kernels and verification of vectorization in compiler IR~\cite{saltyrn,llm-vectorizer}. Extending the guarantee to the complete calling domain requires proofs of the kernels' loops and memory accesses, together with intrinsic semantics that account for numerical differences between their ISAs.

A shared semantic basis for this verification can be constructed at several levels. At the instruction level, ISA specifications describe operations on architectural state, enabling proofs of machine code grounded in detailed hardware semantics~\cite{sail,islaris}. At the level of an intermediate representation, vector operations can be expressed in an IR and compared using its computational and memory semantics. LLVM translation validation and verified instruction selection for Hexagon HVX demonstrate uses of this approach~\cite{alive2,rake}. Intrinsic source code provides a third option: interpreting supported intrinsics as C operations turns the comparison into a C program verification problem.

We choose the C interface to establish functional equivalence under explicit calling contracts. The kernels must agree on their returned outputs and satisfy their memory specifications for admitted inputs and hardware configurations. C source exposes the vector operations alongside the loops and pointers that manipulate them, so these obligations can be expressed without first reconstructing the computation from architectural states. We rely on human experts to check the intrinsic interpretations against ISA documentation, including their numerical and memory behavior. QCP uses these reviewed interpretations as the semantics for the kernel proofs, allowing the surrounding control flow and memory accesses to be verified within the same C framework.

We present \tool, a framework for checking kernel equivalence across ISAs through reusable C semantic shims, with NEON, RVV, AVX, and Hexagon HVX as its target interfaces. \todo{Report the supported intrinsic subsets.} The framework establishes partial correctness relative to the calling contracts and semantic models (Section~\ref{sec:guarantee}). For each kernel pair, \tool binds intrinsic calls to their semantic definitions and states the required output relation in a common specification. QCP symbolically executes the annotated C programs and generates verification conditions for memory accesses, loop invariants, and postconditions, which are discharged through automation and Rocq/Coq proofs~\cite{qcp}. Each implementation is proved against the common specification using its own invariants, so their loops need not advance in lockstep. Implementations, specifications, and proofs remain separate, allowing agents to develop annotations and auxiliary lemmas while retaining the executable kernels.

Human review centers on the calling contract and the evidence produced during verification. An unsuccessful proof attempt can reflect a missing invariant or lemma, so unresolved obligations remain distinct from demonstrated discrepancies. A concrete counterexample must be checked against the original caller's assumptions. If those assumptions exclude the input, the specification can be strengthened to express them, with evidence that all legitimate calls remain covered. If the input is feasible, it identifies a behavioral discrepancy to repair in the target implementation. The revised kernel is then verified again and benchmarked to assess the repair's performance cost.

This equivalence framework supports three uses. First, it checks whether an optimized kernel preserves its reference behavior, within one ISA or across ISAs. Second, it supplies a correctness criterion for porting existing NEON or AVX kernels to RVV or HVX libraries that lack the corresponding operation. Existing implementations provide both source code for translation and a reference for verification. Third, it supports semantic cross-checking of independently developed implementations of the same operation. As in differential testing, disagreement provides a starting point for investigation; equivalence checking makes the relationship between the implementations an explicit proof obligation~\cite{mckeeman,symdiff}. The library contract determines whether a discrepancy violates the required behavior or falls within the permitted numerical differences.

\todo{Evaluation paragraph: experimental evidence and results remain pending.}
We evaluate \tool using the 35 final optimized RVV kernels from SALTyRN, translations of five AVX operations to RVV and HVX, and comparisons of corresponding implementations in XNNPACK~\cite{final-kernels}. We also use \todo{application; candidate: softmax} to examine equivalence across operation boundaries. We verify \todo{N} of the 35 optimized kernel pairs and measure \todo{speedup} relative to the handwritten RVV baseline. The translation and cross-checking studies establish \todo{verified ports} and identify \todo{contract violations, allowed differences, and upstream outcomes}. We report the time and tokens spent constructing proofs, the implementation changes required for repairs, and their effect on performance.

% Alternative short P8, retained at the author's request; not rendered.
% The framework supports validation of kernel optimizations, verified translation
% across ISAs, and cross-checking of existing implementations. The planned main
% study uses SALTyRN's 35 published final optimized RVV kernels. Additional studies
% target AVX-to-RVV/HVX ports, XNNPACK implementation comparisons, and a composed
% SIMD program. All proof counts, repairs, cost measurements, and speedups remain
% \todo{pending evaluation}; none is asserted as a completed contribution.
